\chapter{Accumulation and the periodic law}

\textbf{The figures here are measured on the ideal filling and are deterministic.} The ideal filling
is defined first, because every later claim is a claim about it and not about a measured atom.

\section{The ideal filling}

The element table builds an element's electrons in Madelung order, lowest principal plus azimuthal
first, and Hund's rule inside a subshell, every magnetic orbital taking a spin-up electron before any
takes spin-down. The order fills 1s through 7p, and its capacities run to 118 exactly. It names every
element from hydrogen to oganesson and no further.

This is the necessary condition, the filling the periodic law would follow if nothing competed with
it. Real atoms deviate: chromium, copper and others fill against the order, and their measured
configurations~\cite{src:NIST-Atomic-Spectra-Database} are the reference for a comparison this
research paper does not report. The representation carries the ideal and says so, so that no reader
takes the ideal for a measurement.

\section{Exclusion is an accumulation}

Two fermions may not occupy one state. An electron added to an atom therefore cannot join the
electrons already in the lowest state; it takes the next state the exclusion leaves open, and the one
after that takes the next again. Electrons accumulate into successive states, and the stack of
occupied states is the shell structure. Exclusion is not a bookkeeping rule laid over the atom. It is
the reason the atom has structure at all.

That accumulation is faithful here for the reason the first chapter gives. A reading that rounds two
close states together loses an electron from the stack and shortens a shell. The exact reading never
does, because its discrimination is unbounded, and an accumulation of fermions stays an accumulation
and never collapses. Exactness lets the method hold the stack, and exclusion makes the stack tall.

\section{The recurrence, against a drawn null}

The differentiating electron of an element is the last one the order adds, and its group signature is
its azimuthal type and its subshell's population with the principal number dropped. Dropping the
principal number leaves sodium and lithium reading the same signature, and that recurrence down the
table is the periodic law. It is measured by an equality match rate: how often a signature equals the
one a fixed number of elements ahead.

A match rate means nothing without the rate a shuffle of the same signatures reaches. Two
backgrounds are drawn, each deleting a property from the data. The first keeps the census and deletes
the order: a uniform permutation of the same signatures, which fixes every count exactly. Its floor
is the real floor, and the match rate at the row lengths falls to it. The recurrence is in the order
the shells fill and not in the counts. The second deletes the exclusion: every electron drops to 1s,
an element becomes a count of electrons in one state, every signature is distinct, and the match rate
and its floor are both zero. Without exclusion there is a ladder of counts and no table. The real
accumulation is the positive control, and it is the only arm that carries the recurrence.

\section{The row lengths, off the difference set}

The recurrence does not sit at one lag. The match rate stands above its floor at 32, at 18 and at 8,
and at their sums, because the rows of the table are 2, 8, 8, 18, 18, 32, 32 and not one length. A
reader reporting a single period is wrong for most of the table, the harmonic trap the crystallography
and game theory research papers~\cite{src:Quigg-crystallography,src:Quigg-game-theory} meet from the
other side. What is one number per row is the gap between shell closures, read off the signatures: a
closure is a filled p subshell, with helium closing the first row.

\begin{center}
\begin{tabular}{@{}lrrrrrrr@{}}
\toprule
closure at $Z$ & 2 & 10 & 18 & 36 & 54 & 86 & 118 \\
\midrule
row length     &   & 8  & 8  & 18 & 18 & 32 & 32 \\
\bottomrule
\end{tabular}
\end{center}

\noindent The closures are the noble gases, and the gaps between them are the row lengths, with no
row told to the reader. This is the difference-set reading of the crystallography research paper
moved to one dimension: a period is a difference between two things that agree, and the differences
are the complete candidate set.
